\section{Ablations and Analysis}
\label{sec:ablations}

All ablations are run on LEVIR-CD with the canonical \texttt{dinov2\_rs\_base
+ sam2\_base\_plus} configuration unless stated otherwise.

\subsection{Component ablations}

Table~\ref{tab:ablation-components} removes one GALASH component at a time.

\begin{table}[h]
  \centering
  \small
  \caption{Component ablations on LEVIR-CD test (F1 / IoU).}
  \label{tab:ablation-components}
  \begin{tabular}{lcc}
    \toprule
    Configuration & F1 & IoU \\
    \midrule
    GALASH (full)                            & TODO & TODO \\
    \quad w/o Cross-Change Attention (naive feature diff) & TODO & TODO \\
    \quad w/o latent change loss ($w_\text{lat}{=}0$)     & TODO & TODO \\
    \quad w/o OHEM in BCE                                 & TODO & TODO \\
    \quad w/o Bridge v2 transformer refinement            & TODO & TODO \\
    \quad w/o high-resolution feature tap                 & TODO & TODO \\
    \quad w/o TTA at test time                            & TODO & TODO \\
    \quad SAM decoder fine-tuned ($0.1\times$ LR)         & TODO & TODO \\
    \bottomrule
  \end{tabular}
\end{table}

\paragraph{Cross-Change Attention is essential.} Replacing it with naive
feature differencing ($Z^{(c)} = Z^{(r)} - Z^{(t)}$) drops F1 by
\textsc{TODO} points, confirming that the attention pattern --- not the
feature difference --- is the useful signal. The learned temperature
$\tau = e^\lambda$ converges to approximately \textsc{TODO} after training.

\paragraph{Latent change loss acts as a regulariser.} Removing
$\mathcal{L}_\text{lat}$ produces a small drop in F1 (\textsc{TODO} points)
and causes a visible reduction in cross-dataset generalisation
(Table~\ref{tab:xdataset}), consistent with its role as a shape-agnostic,
patch-level supervision signal.

\paragraph{SAM decoder fine-tuning.} Unfreezing SAM's mask decoder at
0.1$\times$ the main learning rate yields a modest F1 improvement on
LEVIR-CD but costs \textsc{TODO}\% more trainable parameters and increases
the risk of catastrophic forgetting on zero-shot datasets.

\subsection{Parameter and FLOP efficiency}

Table~\ref{tab:params} compares GALASH to recent methods in terms of
trainable parameter count. All numbers are for the default LEVIR-CD
configuration.

\begin{table}[h]
  \centering
  \small
  \caption{Trainable parameter counts vs.\ F1 on LEVIR-CD. GALASH trains
    16.5~M parameters --- fewer than ChangeFormer, DDPM-CD, and SChanger,
    and competitive with BIT.}
  \label{tab:params}
  \begin{tabular}{lcc}
    \toprule
    Method & Trainable params (M) & F1 \\
    \midrule
    TinyCD~\citep{codegoni2023tinycd}       & 0.28  & 91.05 \\
    BIT~\citep{chen2021remote}              & 3.00  & 89.31 \\
    FC-Siam-diff~\citep{daudt2018fully}     & 1.35  & 86.31 \\
    BAN-BiT~\citep{li2024ban}               & 3.80  & 91.85 \\
    ChangeFormer~\citep{bandara2022changeformer} & 41.0  & 91.11 \\
    DDPM-CD~\citep{bandara2024ddpmcd}       & 45.7  & 90.91 \\
    \midrule
    GALASH (ours)                           & \textbf{16.5} & \textbf{TODO} \\
    \bottomrule
  \end{tabular}
\end{table}

\subsection{Qualitative results}

Figure~\ref{fig:qualitative} (TODO) shows change maps predicted by GALASH on
challenging examples from each benchmark. We highlight three failure modes
that motivate directions for future work: (i) very small building additions
in rural S2Looking scenes, (ii) seasonal colour shifts misclassified as
changes on CDD, and (iii) registration-induced edge artefacts on SECOND.
